% Methods / position paper: a hierarchical Bayesian interaction model as a
% standard for analyzing LLM evaluations. Case study reuses the retention data.
\documentclass{llncs}

\usepackage[T1]{fontenc}
\usepackage{graphicx}
\usepackage{amsmath}
\usepackage{amssymb}
\usepackage{booktabs}
\usepackage{microtype}
\usepackage[htt]{hyphenat}
\emergencystretch=1.5em
\graphicspath{{figures/}{../report-figures/v6/}}
\usepackage[style=apa]{biblatex}
\addbibresource{references.bib}

\begin{document}

\title{Evaluations Are Crossed Designs: A Hierarchical Bayesian
Interaction Model for Analyzing Model Evaluations}
\titlerunning{Evaluations Are Crossed Designs}
\author{Tom\'as P. Korenblit\orcidID{0009-0002-5682-8475}}
\authorrunning{T. P. Korenblit}
\institute{Universidad Nacional de San Mart\'in, Buenos Aires, Argentina\\
\email{tomaskorenblit@gmail.com}}
\maketitle

\begin{abstract}
Most language-model evaluations report a mean score, sometimes with a
standard error. This discards three things the data contains: the scores
are usually graded rather than binary, the design is models crossed with
items, and models have item-specific strengths, a model-by-item
interaction. We argue that the appropriate analysis is a hierarchical
Bayesian graded-response model with crossed random effects for models and
items and an explicit model-by-item interaction term, fit by MCMC. This sits
in the explanatory item-response and generalizability-theory traditions but
fills a gap: existing item-response work in machine learning is almost
entirely binary, variational, and aimed at subselecting items rather than
analyzing the systems under test. We demonstrate the model on a
multi-turn instruction-retention evaluation, where a ladder of models shows
the interaction term improving out-of-sample fit by about 59 ELPD (five
standard errors) over additive specifications, while additive per-model
terms add nothing. Reporting raw means or additive effects on this data
misattributes interaction variance to main effects and manufactures a
spurious finding that the interaction model dissolves. We give a concrete
specification, priors, and a reporting checklist, and discuss when the model
is and is not worth its cost.
\keywords{model evaluation \and item response theory \and Bayesian
inference \and generalizability theory \and measurement}
\end{abstract}

\section{Introduction}
\label{sec:intro}

A language-model evaluation produces a table: systems in rows, test items in
columns, a score in each cell. The near-universal summary is a column mean
per system, occasionally with a standard error \parencite{miller2024errorbars}.
This summary throws away most of the structure in the table. First, the
scores are frequently \emph{graded}: a rubric returns 0--3, a judge returns a
Likert rating, a partial-credit checker returns degrees of compliance.
Averaging treats these as interval-scaled, which they are not. Second, the
design is \emph{crossed}: every system meets every item, so item difficulty
and system ability are confounded in any per-system mean, and a system that
draws an easier item subset looks better for free. Third, and most
consequentially, systems have \emph{item-specific strengths}: a model strong
on one task may be weak on another in a way not predicted by its average.
This is a model-by-item interaction, and a mean by construction cannot see
it.

We argue for a specific analysis as a default: a hierarchical Bayesian
graded-response model with crossed random effects for systems and items and
an explicit model-by-item interaction term, estimated by MCMC. The
ingredients are standard in psychometrics, item response theory (IRT)
\parencite{samejima1969grm, deboeck2004explanatory}, crossed-effects modeling
\parencite{baayen2008crossed}, and generalizability theory
\parencite{brennan2001gtheory}, and are well supported by modern tooling
\parencite{burkner2021irt, mcelreath2020rethinking, pymc2023}. What is new is
assembling them for the evaluation setting: treating evaluation scores as
ordinal, putting the experimental manipulation as a covariate on ability, and
making the model-by-item interaction a first-class, identifiable term.

Our contributions are: (i) we characterize evaluation data as a crossed,
graded, interaction-bearing design and state the model that matches it;
(ii) we show, on a real evaluation, that the interaction term is the dominant
structure, $\approx$59 ELPD better than additive models, and that ignoring
it (raw means or additive effects) produces a concrete, published error;
(iii) we give a reproducible specification, including why an interaction
variance component is preferable to a low-rank factor model on
identifiability grounds; and (iv) we provide a reporting checklist and an
honest account of the method's costs.

\section{What the field does, and the gap}
\label{sec:gap}

Three lines of work move beyond the raw mean, each addressing one facet but
not the whole.

\paragraph{Uncertainty.} \textcite{miller2024errorbars} treats an eval as an
experiment and imports standard errors, clustering, and paired differences;
\textcite{bowyer2025clt} shows the normal approximation underestimates
uncertainty at the small sample sizes typical of specialized evals. This
fixes the error bar on a mean but keeps the mean: no item structure, no
interaction.

\paragraph{Item response theory.} IRT has been used to build evaluation
scales \parencite{lalor2016irt}, reweight leaderboards by item
informativeness \parencite{rodriguez2021leaderboards}, compare test sets
\parencite{vania2021testsets}, and construct efficient benchmarks
\parencite{polo2024tinybenchmarks}. This literature is almost entirely
\emph{binary} (correct/incorrect), estimated by \emph{variational}
approximation for scale, and aimed at \emph{item subselection} rather than
analyzing the systems under test. The one graded-response application we
found diagnoses judge reliability, not the systems being evaluated
\parencite{choi2026judgeirt}.

\paragraph{Measurement and ranking.} Construct-validity work asks whether
benchmarks measure what they claim \parencite{bean2025construct};
generalizability theory decomposes score variance across facets but in an
ANOVA, non-generative form \parencite{song2025gtheory}; preference ranking
uses Bradley--Terry with intervals \parencite{chiang2024arena} but for
pairwise outcomes, not graded item scores. LLM-as-judge reliability is
typically reported as raw agreement, which is not chance-corrected
\parencite{zheng2023judge}.

The gap is the intersection: no standard approach simultaneously (a) treats
scores as \emph{ordinal}, (b) models systems and items as \emph{crossed with
partial pooling}, and (c) includes an explicit \emph{model-by-item
interaction}. The pieces exist in adjacent literatures; assembling them is
the contribution.

\section{The recommended model}
\label{sec:model}

Let $y_{i}\in\{0,\dots,K\}$ be the graded score for observation $i$ with
system $m[i]$, item $d[i]$, and an optional experimental covariate $t_i$ (a
condition, intervention, or treatment). The model is a cumulative-link
(ordered-logistic) graded-response model,
\[
\Pr(y_i \le k) = \mathrm{logistic}(\kappa_k - \eta_i),\qquad
\kappa_1 < \dots < \kappa_K,
\]
with linear predictor
\[
\eta_i = \alpha_{d[i]} + \beta_{d[i]}\, t_i
       + g_{m[i]} + s_{m[i]}\, t_i
       + z_{m[i],d[i]}\, t_i .
\]
Here $\alpha_d$ is per-item baseline difficulty, $\beta_d$ the per-item
effect of the covariate, $g_m$ the per-system ability, $s_m$ the per-system
covariate sensitivity, and $z_{m,d}\sim\mathcal{N}(0,\sigma_{\text{int}})$
the \textbf{model-by-item interaction}. All group effects are non-centered;
cutpoints use an offset parameterization; SDs take weakly-informative
half-normal/exponential priors and correlations, where used, an LKJ prior
\parencite{mcelreath2020rethinking, burkner2021irt}.

Four choices deserve justification. \emph{Ordinal likelihood}: graded scores
are not interval-scaled, so a cumulative link is the principled choice over a
Gaussian or a binarization. \emph{Crossed effects}: systems and items are
fully crossed, so crossed random intercepts remove the confound between item
difficulty and system ability that a per-system mean carries; this also
yields a generalizability-theory variance decomposition for free
\parencite{brennan2001gtheory}. \emph{Interaction as a variance component}:
the non-additivity can be written as a low-rank product $\lambda_m\varphi_d$
(a factor / 2PL form), but that is unidentified up to sign and scale and in
practice multimodal; the additive cell effect $z_{m,d}\sim\mathcal{N}(0,
\sigma_{\text{int}})$ is fully identified, samples cleanly, and gives a
single interpretable estimand, the interaction magnitude. \emph{Full
Bayesian estimation}: MCMC propagates uncertainty into every downstream
quantity, including system rankings, where overlapping intervals correctly
read as ``tied''; partial pooling shrinks noisy per-cell estimates, the
behavior that makes small-$n$ evals tractable.

What the model returns that a mean does not: per-item and per-system effects
with calibrated intervals; a variance decomposition (how much of the
variation is item, system, interaction); model comparison by leave-one-out
cross-validation \parencite{vehtari2017loo} to test whether structure is
warranted; and, for judge-scored evals, a place to put rater effects and to
report chance-corrected agreement rather than raw agreement.

\section{Case study: the interaction is the structure}
\label{sec:case}

We apply the framework to a multi-turn instruction-retention evaluation
\parencite{korenblit2026}: five models are scored on whether they retain
twelve coding instructions across three codebases, under a baseline and a
treatment (stating the instruction), on a 0--3 compliance scale, 1{,}067
observations. We fit a ladder of increasingly structured models and compare
them by leave-one-out cross-validation (Table~\ref{tab:loo},
Fig.~\ref{fig:loo}).

\begin{table}[t]
\centering
\caption{Leave-one-out comparison on the retention evaluation. Models with a
model-by-item interaction (M8, M\_irt) fit far better than additive models.}
\label{tab:loo}
\begin{tabular}{llrr}
\toprule
Model & Structure & ELPD & $\Delta$ELPD (SE) \\
\midrule
M8 & + model$\times$item interaction (variance) & $-545.9$ & --- \\
M\_irt & + 2PL discrimination $a_d\theta_m$ & $-555.4$ & $9.5\ (10.5)$ \\
M\_corr & additive + correlated slopes & $-604.4$ & $58.5\ (11.2)$ \\
M4 & additive per-model + per-item & $-604.8$ & $58.9\ (11.2)$ \\
M0 & per-item only (no system structure) & $-607.2$ & $61.3\ (11.7)$ \\
\bottomrule
\end{tabular}
\end{table}

\begin{figure}[t]
\centering
\includegraphics[width=0.85\textwidth]{fig_ladder_loo}
\caption{Out-of-sample fit (ELPD) with standard errors. The two models that
include a model-by-item interaction (teal) separate cleanly from the
additive models (grey), which are indistinguishable from each other.}
\label{fig:loo}
\end{figure}

The result is unambiguous. The two specifications that capture a
model-by-item interaction fit decisively better than every additive model,
by about 59 ELPD, more than five standard errors. The two interaction forms
(an interaction variance component, and an IRT discrimination term) are
indistinguishable from each other, so the interaction matters, not its exact
parameterization. The additive embellishments do nothing: a per-system main
effect (M4) is indistinguishable from a model with no system structure at
all (M0), and correlating system intercepts with slopes (M\_corr) does not
help. Under M8 the interaction is the largest structured source of variance,
$\sigma_{\text{int}}=1.84$ (94\% HDI $[1.26, 2.40]$), exceeding the per-item,
per-system, and per-system-slope components.

This is not a curiosity; it changes published conclusions. The mean and the
additive model attribute the interaction variance to main effects. In this
dataset that inflated the per-item effect (its SD nearly halves when the
interaction is admitted) and, combined with a scoring bug, produced a
reported finding that one instruction was \emph{harmed} by being stated. The
interaction model localizes that effect to a single extreme (system, item)
cell of one model, and a separate framing experiment confirms it is an
artifact, not a real effect. An analysis that cannot represent the
interaction cannot diagnose this.

\section{A reporting checklist}
\label{sec:checklist}

For an evaluation of systems crossed with graded items we recommend:
(1) model scores with an ordinal/cumulative link, not a mean over a graded
scale; (2) include crossed random effects for systems and items with partial
pooling; (3) include a model-by-item interaction term, and report
$\sigma_{\text{int}}$ alongside the main-effect SDs; (4) compare
specifications by leave-one-out cross-validation rather than asserting a
structure; (5) report rankings with their posterior intervals, treating
overlap as ``tied''; (6) check prior sensitivity of the variance components,
which are prior-influenced when the number of systems or items is small; (7)
for judge-scored evals, model rater effects and report chance-corrected
agreement. Items (1)--(3) are the substance; (4)--(7) are the discipline that
keeps the conclusions honest.

\section{Limitations and when not to use this}
\label{sec:limits}

The variance components are prior-sensitive when groups are few; with a
handful of systems or items, report rankings and the qualitative
interaction conclusion, not precise variance magnitudes. MCMC is more
expensive than the variational estimation common in IRT-for-eval work; for
very large item banks aimed at subselection, amortized variational methods
remain appropriate, and the model here is better suited to analyzing a fixed
set of systems than to scaling to millions of items. When scores are
genuinely binary and items are not crossed with systems in a way that admits
interaction, the full model reduces to simpler, well-understood special
cases, and the added structure is not worth its cost. The recommendation is
for the common middle case: a moderate number of systems, graded items, a
crossed design, and a comparative question.

\section{Conclusion}
Evaluation tables are crossed designs of graded items, and the quantity that
matters is often how systems differ item by item, an interaction a mean
cannot represent. A hierarchical Bayesian graded-response model with crossed
effects and a model-by-item interaction matches this structure, is supported
by established psychometric theory and modern tooling, and, on a real
evaluation, reveals structure that additive analyses miss by 59 ELPD and that
raw means turn into a spurious finding. We recommend it as a default for
evaluation analysis and provide the specification and checklist to apply it.

\subsubsection*{Use of large language models.}
During the preparation of this work the author used a large language model
(Claude) to improve the language and readability of the manuscript and to
assist with analysis and figure code; all results were computed by the
statistical models and verified by the author.

\printbibliography

\end{document}
